\documentclass[12pt, a4paper]{article}
\usepackage[T1]{fontenc}
\usepackage{amsmath, amssymb}
\usepackage{geometry}
\geometry{margin=2.5cm}
\usepackage{xcolor}
\usepackage{hyperref}
\hypersetup{colorlinks=true, urlcolor=blue!60!black, linkcolor=black}
\usepackage{listings}
\usepackage{booktabs}
\usepackage{tabularx}
\usepackage{xeCJK}
\setCJKmainfont{PingFang TC}

\lstset{
  basicstyle=\ttfamily\small,
  breaklines=true,
  frame=single,
  backgroundcolor=\color{gray!10}
}

\newcolumntype{Y}{>{\raggedright\arraybackslash}X}

\title{我的 AI Agent 使用手冊：把工具清單變成工作流}
\author{Sheng-Yuan Chen}
\date{2026-06-11}

\begin{document}
\maketitle

\section{前言：為什麼要寫一本「使用手冊」}

用 Claude Code 一段時間後，我累積了一大堆 skills——研究、寫作、投影片、數據分析、發筆記，零零總總幾十個。問題是：\textbf{工具一多，反而不知道什麼時候該用哪個}。更常見的狀況是，同一個 session 裡東問西問，最後 AI 開始「忘東忘西」、品質下滑。

於是我把自己的用法整理成一本手冊，回答三個問題：（1）我有哪些工具、什麼情況用哪個？（2）怎麼把零散的工具串成「工作流」，而不是每次重想？（3）怎麼有效率地用 AI agent——省 token、管理 context、寫好 \texttt{CLAUDE.md}、什麼任務該開哪家 AI？這篇是那本手冊的公開版，內容裡所有關於 Claude Code 的事實都對照了 Anthropic 官方文件。

\section{先建立正確的心智模型}

\subsection{AI agent $\neq$ 聊天機器人}

聊天機器人是「你問、它答、它等你」。\textbf{Agent（像 Claude Code）是「你描述目標，它自己讀檔、跑指令、改檔、驗證、迭代」}——你在旁邊看、必要時導正，或直接走開。對研究的意義是：別再「自己想好每一步、一句句下指令」，而是把目標、限制、驗收標準講清楚，讓它自己跑。比方說，不要說「幫我把這個變數 demean」，而是說：「用 first-difference 估這個 panel DML，跑完給我 bias / SE 表，並和 OLS 對照。」

\subsection{三個延伸機制：Skills / CLAUDE.md / Memory}

\begin{table}[h]
\centering
\begin{tabularx}{\linewidth}{lYY}
\toprule
機制 & 是什麼 & 角色 \\
\midrule
Skills & 一包 \texttt{SKILL.md} ＋ 腳本，給 AI 特定領域知識與可重複流程 & 只在相關時才載入的「專門流程」 \\
CLAUDE.md & 每次開場自動載入的「專案說明書」 & 每次都載入的「常駐規則」 \\
Memory & 跨 session 的長期記憶（偏好、專案狀態） & 你的「個人筆記」 \\
\bottomrule
\end{tabularx}
\end{table}

一句話記住分工：\textbf{CLAUDE.md ＝每次都載入的常駐說明；Skill ＝需要時才載入的專門流程；Memory ＝你的偏好與狀態}。三者別混用——這在省 token 的部分很關鍵。

\section{把技能整理成「工作流」，而不是「工具清單」}

工具清單會讓你每次都重想；\textbf{工作流}則是把「某個常見任務」固定成一條串接。以下是幾個對研究生最實用的例子（觸發句直接貼給 agent 即可）。

\begin{itemize}
  \item \textbf{讀懂並報告一篇論文：}讀 PDF $\rightarrow$ 建立「理解 $\rightarrow$ 大綱 $\rightarrow$ 逐張 frame $\rightarrow$ 版面」的中介表示 $\rightarrow$ 寫 \texttt{.tex} $\rightarrow$ 逐頁視覺檢查。別讓 AI 從 PDF 直接跳到投影片，中間要有可檢查、可修補的階段。
  \item \textbf{多篇文獻找研究缺口：}做跨論文 synthesis（不要逐篇摘要），把 gap 分類並評 Strong / Medium / Weak，每個 claim 標出論文出處，不確定的標 \texttt{to verify}。
  \item \textbf{從想法到論文初稿：}先讓 AI 用問答把想法逼成一份 \texttt{SPEC.md}（它會問你沒想到的取捨與 edge case），再開新 session 執行——新 session context 乾淨、又有 spec 可循。
  \item \textbf{數據分析到出表／圖：}要求可重現的腳本 ＋ 證據；\texttt{data/raw/} 不要動，清理後存 \texttt{data/processed/}。
\end{itemize}

\textbf{共通原則：給它一個能跑的「驗收檢查」。}最有效的一招是給 AI 一個能自己跑、會回傳 pass / fail 的檢查（測試、編譯、跑回歸對照、screenshot 比對）。有了它，AI 會自己跑到通過，你不用當「人肉驗證迴圈」；並且要求它給證據，而不是一句「我做好了」。

\section{高效使用 AI agent 的核心：context 是最稀缺的資源}

\subsection{為什麼 AI 用久了會「變笨」：context rot}

AI 的 \textbf{context window（脈絡視窗）}裝著整段對話——每則訊息、每個被讀的檔、每次指令輸出。它會很快填滿，而且越滿、表現越差：模型開始「忘記」前面的指令、變多錯。Anthropic 把這現象叫 context rot：transformer 要處理 $n^2$ 的 token 兩兩關係，context 越長、注意力被攤越薄。結果是準確率隨 context 變長而漸降（是斜坡、不是斷崖）。核心心法一句話：\textbf{用「最小的高訊號 token 集合」達成任務}，不是塞越多越好。

\subsection{省 token 的具體做法}

\begin{enumerate}
  \item \textbf{任務之間 \texttt{/clear}。}換到不相關的工作就清空 context；最常見的浪費就是「kitchen sink session」。
  \item \textbf{同一問題被你糾正超過兩次 $\rightarrow$ \texttt{/clear} 重開}，用更好的 prompt。乾淨 session ＋ 好 prompt 幾乎總是贏過拖很長的 session。
  \item \textbf{階段做完就 \texttt{/compact}}（不是等到變笨才壓），可加指示如 \texttt{/compact 著重在 R 程式}。
  \item \textbf{用 \texttt{/context} 看誰在吃 context}。
  \item \textbf{要 AI 去讀檔，而不是貼長文進來}。用 \texttt{@檔名} 引用、或叫它自己讀檔／抓網址。
  \item \textbf{大量讀檔的探索丟給 subagent}——它在獨立 context 探索，只回報約 1{,}000--2{,}000 token 的摘要。
  \item \textbf{小問題用 side question}，答案不進對話歷史。
  \item \textbf{\texttt{CLAUDE.md} 保持精簡}——它每次都載入、整段常駐。
\end{enumerate}

\subsection{CLAUDE.md 是什麼、怎麼寫}

定義：每次對話開場 AI 都會自動讀的「專案說明書」。放它從程式碼推不出來、但每次都需要的東西：常用指令、程式風格、工作流規則、專案慣例。位置會自動合併：

\begin{table}[h]
\centering
\begin{tabularx}{\linewidth}{lY}
\toprule
位置 & 作用 \\
\midrule
\texttt{\textasciitilde/.claude/CLAUDE.md} & 全域，套用到所有 session \\
\texttt{./CLAUDE.md}（專案根） & 專案層，commit 進 git 與人共享 \\
\texttt{./CLAUDE.local.md} & 個人專案筆記，記得加進 \texttt{.gitignore} \\
父／子目錄 & monorepo 自動往上拉、讀子目錄檔時才拉子目錄的 \\
\bottomrule
\end{tabularx}
\end{table}

官方寫法守則：保持短、人類可讀，每一行問「拿掉這行會讓 AI 犯錯嗎？」不會就刪；\textbf{臃腫的 \texttt{CLAUDE.md} 會讓 AI 忽略你真正的指令}。該放：猜不到的指令、與預設不同的風格、測試方式、架構決定、環境怪癖。不該放：讀程式就知道的事、會頻繁變動的資訊、長篇教學。加 \texttt{IMPORTANT} / \texttt{YOU MUST} 可提高遵從度；在對話裡按 \texttt{\#} 可把一句話直接寫進 \texttt{CLAUDE.md}。我自己把專案根的 \texttt{CLAUDE.md} 當成路由表，只寫「哪個任務先讀哪個 \texttt{STATUS.md}」，把會變動的進度交給各專案的 \texttt{STATUS.md}。

\subsection{先探索 $\rightarrow$ 再計畫 $\rightarrow$ 才動手}

官方建議四階段：Explore（plan mode 只讀不改）$\rightarrow$ Plan $\rightarrow$ Implement $\rightarrow$ Commit。但 plan mode 有 overhead：小事直接做，只在「不確定做法、會動多個檔、對該段不熟」時才用。一句話能描述改動，就跳過計畫。

\section{不同 AI 的優勢與選用}

模型版本與 benchmark 變很快，下面是 2026 年中的「相對強項」歸納，確切數字請臨用前再查。

\begin{table}[h]
\centering
\begin{tabularx}{\linewidth}{lYY}
\toprule
模型家族 & 相對強項 & 什麼情況開它 \\
\midrule
Claude（Anthropic） & 程式／agentic 任務、複雜邏輯與除錯、長流程、最不「AI 腔」的英文寫作 & agentic 工作流、寫改 R／LaTeX、用自己口吻寫英文論文 \\
GPT（OpenAI） & 通用、知識面廣、結構化推理、外掛生態 & 快速廣博的解法、brainstorming、工具整合 \\
Gemini（Google） & 速度快、context window 超大、原生 Google 搜尋、多模態 & 一次塞超大量文字、即時網路資訊、處理圖片影片 \\
\bottomrule
\end{tabularx}
\end{table}

實務建議：agentic 流程 ＋ 寫 code ＋ 英文論文寫作走 Claude / Claude Code；超長一次性閱讀可考慮 Gemini 的大 context；快速查證用帶搜尋的模型。不要迷信單一 benchmark 冠軍——挑模型看任務類型。Claude 家族內部：Opus 做最難的推理與長流程、Sonnet 是又快又強的日常主力、Haiku 處理量大重複的小事。

\section{小結}

\begin{table}[h]
\centering
\begin{tabularx}{\linewidth}{lY}
\toprule
主題 & 一句話心法 \\
\midrule
心智模型 & 描述目標與驗收標準，讓 agent 自己跑，別當它的手腳 \\
工作流 & 把常見任務固定成串接，附「觸發句」，每次複製即用 \\
驗收 & 給一個能自己跑、回傳 pass/fail 的檢查；要證據不要「我做好了」 \\
context & 它是最稀缺資源；越滿越笨（context rot） \\
省 token & 任務間 \texttt{/clear}、階段後 \texttt{/compact}、讀檔別貼長文、探索丟 subagent \\
CLAUDE.md & 短、只放猜不到又每次要的；會變動的進度交給 \texttt{STATUS.md} \\
選 AI & 看任務類型選家族；Claude 做 agentic／寫作、Gemini 大 context、GPT 廣博 \\
\bottomrule
\end{tabularx}
\end{table}

工具會一直變多，但「先理解、再規劃、才動手；管好 context；給它能驗證的檢查」這幾條原則不太會過期。

\vspace{1em}
\noindent\textbf{參考來源：} Anthropic, \emph{Best practices for Claude Code}；\emph{Manage Claude's memory}；\emph{Effective context engineering for AI agents}。模型比較為第三方 2026 年資料，數字會過期，僅供方向參考。

\end{document}
